\section*{Chapter \ref{ch:m2:the-fx-market} --- The FX Market}

\begin{solution}{exo:m2:the-fx-market:1}
It sells at the dealer's bid: $5\,000\,000 \times 1.1462 = \text{USD}~5\,731\,000$.
The spread is 2 \omterm{def:m2:the-fx-market:pair}{pips}, 1.74 basis points of the mid.
\end{solution}

\begin{solution}{exo:m2:the-fx-market:2}
USDJPY, \omterm{def:m2:the-fx-market:pair}{pip} 0.01; EURGBP, \omterm{def:m2:the-fx-market:pair}{pip} 0.0001; AUDJPY, \omterm{def:m2:the-fx-market:pair}{pip} 0.01.
\end{solution}

\begin{solution}{exo:m2:the-fx-market:3}
Thursday 8 October with no holiday; Friday 9 October with a Japanese holiday on
Wednesday; Friday 9 October with a US holiday on Thursday, the value date itself.
USDCAD settles on Wednesday 7 October.
\end{solution}

\begin{solution}{exo:m2:the-fx-market:4}
Bid $1.1462/1.3373 = 0.85710$, ask $1.1464/1.3371 = 0.85738$: 2.78 \omterm{def:m2:the-fx-market:pair}{pips}.
\end{solution}

\begin{solution}{exo:m2:the-fx-market:5}
Each trade involves two currencies and each is counted, so the shares of all
currencies sum to twice the total. With the dollar on one side of 89\% of trades,
a cross such as SEKJPY is thin: most of it is traded as two dollar trades, SEK
against USD and USD against JPY, where liquidity is deep.
\end{solution}

\begin{solution}{exo:m2:the-fx-market:6}
The trade is given up to the prime broker, which becomes the executing dealer's
counterparty; the dealer faces the prime broker's credit, not the fund's. The
prime broker faces the fund on an offsetting trade and bears its credit risk,
which it limits and margins.
\end{solution}

\begin{solution}{exo:m2:the-fx-market:7}
209.7375/209.7956, a relative spread of 2.771 basis points; the legs'
relative spreads are 1.496 and 1.275, which sum to 2.771.
\end{solution}

\begin{solution}{exo:m2:the-fx-market:8}
The total is spread over hundreds of pairs, instruments and time zones; most of
it is FX swaps and dealer-to-dealer trading, and a small currency against another
may trade a few million an hour. Liquidity at the screen price is only the
size shown; a large order is worked, pays a wider spread or moves the price,
and at some times of day and in some pairs there is little to trade at all.
\end{solution}

\begin{solution}{pb:m2:the-fx-market:1}
\textbf{1.} Both asks: it buys dollars with yen at the USDJPY ask, 156.88, and
euros with dollars at the EURUSD ask, 1.1464.
\textbf{2.} 179.7929/179.8472.
\textbf{3.} 5.43 \omterm{def:m2:the-fx-market:pair}{pips}, 3.02 basis points.
\textbf{4.} 1.745 and 1.275 basis points, summing to 3.02: relative spreads add.
\textbf{5.} USD~11\,464\,000, bought and immediately spent.
\textbf{6.} Yes: its bid, 179.80, is below the synthetic ask and its ask, 179.84,
above the synthetic bid; it is also tighter than the synthetic.
\textbf{7.} No: its bid, 179.85, is above the synthetic ask. Sell euros to B at
179.85 and buy them back through the legs at 179.8472.
\textbf{8.} $10\,000\,000 \times (179.85 - 179.84723)$, JPY~27\,680, about USD~176.
\textbf{9.} A quote left behind a move in the legs, by a slow system or a
distracted trader; arbitrageurs and the dealer's own checks remove it within
moments.
\textbf{10.} From dealer~A at 179.84, below both B's ask and the synthetic ask.
\textbf{11.} Two tickets and two spreads, the risk that the second leg moves
before it is done, three currencies to settle, and possibly minimum sizes on the
legs.
\textbf{12.} When a dealer has offsetting cross flow or sees two-way interest in
the cross itself, and can quote inside the synthetic.
\textbf{13.} The direct trade exchanges euros and yen; the synthetic one also
buys and sells dollars, which must be delivered and paid, usually through CLS.
\textbf{14.} The yen legs and the direct cross settle on Friday 9 October; the
EURUSD leg on Thursday 8 October. The dollars bought on Thursday are needed on
Friday: the synthetic route needs a one-day swap to line them up.
\textbf{15.} It nets the cross against other clients' flows and hedges only the
remainder, often at better prices than a client trading the legs.
\textbf{16.} If the legs move and the quote does not, it is crossed with the
synthetic and is picked off, as dealer~B's was.
\textbf{17.} The dollar pairs are the deep markets; a thin cross is quoted from
them and its own book adds little.
\textbf{18.} Buying EURJPY in size through the legs buys EURUSD and USDJPY:
both rise, and so does the synthetic cross.
\textbf{19.} Named result: \emph{the synthetic bound} for EURJPY is
179.7929/179.8472, 5.43 \omterm{def:m2:the-fx-market:pair}{pips}: the client should never pay more than 179.8472 for
euros, since it can buy them through the legs; any direct quote is a gain only if
it is inside that.
\textbf{20.} The bid and ask of the two dollar pairs, through which anyone can
trade it.
\end{solution}

\begin{solution}{iq:m2:the-fx-market:1}
The dealer buys euros at 1.1462 dollars and sells them at 1.1464; the spread is
two \omterm{def:m2:the-fx-market:pair}{pips}. A client who buys euros hits the ask, paying 1.1464 dollars for each
euro.
\iqlookfor{base and quote, and the side a buyer pays.}
\end{solution}

\begin{solution}{iq:m2:the-fx-market:2}
Take the trade date (after the 17:00 New York roll), move to the second day that
is a business day in both currencies, ignoring a US holiday on the intermediate
day for dollar pairs, and to the first for USDCAD; then check the value date is
open in both. Test weekends, holidays on the intermediate and on the value day
in each currency, US holidays in the middle, year ends, trades just before and
after 17:00, and the T+1 pairs.
\iqlookfor{the rule, its dollar exception, and a real list of edge cases.}
\end{solution}

\begin{solution}{iq:m2:the-fx-market:3}
Selling euros for pounds through dollars: sell euros at the EURUSD bid $b_1$,
buy pounds with those dollars at the GBPUSD ask $a_3$: bid $b_1/a_3$. Buying
euros: pay the EURUSD ask $a_1$ with dollars from selling pounds at the GBPUSD
bid $b_3$: ask $a_1/b_3$.
\iqlookfor{the correct sides of each leg.}
\end{solution}

\begin{solution}{iq:m2:the-fx-market:4}
Its participants are banks that trade with each other and their clients in
large sizes, across time zones, in many currencies, and settle bilaterally;
no single exchange or clearing house was needed or imposed. Electronic trading
added interdealer order books, then multi-dealer and \omterm{def:m2:the-fx-market:sdp}{single-dealer platforms},
algorithms and non-bank market makers; it lowered costs and fragmented
liquidity, and concentrated dealing in the banks with the most flow.
\iqlookfor{history and the effects of electronification.}
\end{solution}

\begin{solution}{iq:m2:the-fx-market:5}
It lets the client trade with many dealers in the prime broker's name, on its
credit, and settles and reports everything in one place; the dealers face the
prime broker. It takes the client's credit risk: if the client fails, the prime
broker must honour the trades. It limits that risk with trading limits, margin
and intraday monitoring (\cref{ch:m2:getting-access-fx}).
\iqlookfor{the give-up mechanism and the credit risk.}
\end{solution}

\begin{solution}{iq:m2:the-fx-market:6}
Keep each dollar leg's best bid and ask in a fixed array, updated from the market
data handlers; on each leg update recompute only the crosses that use it, from
a precomputed dependency table, with the side rules resolved at start-up so the
hot path is two multiplications or divisions. Publish with a sequence number
and the legs' timestamps, suppress a cross when a leg is stale, and test against
a slow reference implementation on recorded data.
\iqlookfor{precomputed dependencies, no allocation on the hot path, staleness
checks.}
\end{solution}
